% preamble.tex -- shared preamble of the written proofs of the Lean-checked results.
% Every section file is written for this preamble; do not define new packages or macros inside a section file.
\usepackage[margin=0.9in]{geometry}
\usepackage[T1]{fontenc}
\usepackage{lmodern}
\usepackage{amsmath,amssymb,amsthm,mathtools}
\usepackage{booktabs,array,longtable}
\usepackage{enumitem}
\usepackage{microtype}
\usepackage{url}
\usepackage[hidelinks]{hyperref}

\setlist{itemsep=0pt,topsep=2pt,parsep=0pt}
\emergencystretch=2em
\def\UrlBreaks{\do\.\do\/\do\_}

% Statement heads carry long Lean names. amsthm sets the head in a box and unboxes it into the first
% paragraph; a high tolerance inside every statement lets TeX use the break points of the head.
\makeatletter
\def\es@sloppy{\tolerance 9999 \hfuzz .5\p@ \ifdim\emergencystretch<3em \emergencystretch 3em\fi}
\g@addto@macro\th@plain{\es@sloppy}
\g@addto@macro\th@definition{\es@sloppy}
\g@addto@macro\th@remark{\es@sloppy}

% A run-in heading (\subsection, \paragraph) followed directly by a statement, a list or a proof is set as a
% one-line paragraph, after which the list adds \@beginparpenalty (negative), an invitation to break the page
% below the heading. Forbid that break; and until the next paragraph starts, let a \label (or another command
% that ends with \@esphack) in vertical mode add \nobreak after its \write, which would otherwise be a break point.
\newif\ifes@keephead
\let\es@trivlist@orig\@trivlist
\def\@trivlist{\if@noskipsec
    \@beginparpenalty\@M\es@trivlist@orig\global\es@keepheadtrue
  \else\es@trivlist@orig\fi}
\let\es@esphack@orig\@esphack
\def\@esphack{\es@esphack@orig\ifvmode\ifes@keephead\nobreak\fi\fi}
\AddToHook{para/begin}{\global\es@keepheadfalse}

% Paragraph headings: bold and run in. A \paragraph that directly follows a run-in heading keeps that heading
% on its page (no \@secpenalty between them).
\def\paragraph{\if@noskipsec\leavevmode\par\nobreak\@nobreaktrue\fi
  \@startsection{paragraph}{4}\z@\z@{-\fontdimen2\font}{\normalfont\bfseries}}
\makeatother
% Hyphenation exceptions (not macros): the default patterns would break ``namespace'' as ``names-pace''.
\hyphenation{name-space name-spaces}

\newtheorem{theorem}{Theorem}[section]
\newtheorem{lemma}[theorem]{Lemma}
\newtheorem{proposition}[theorem]{Proposition}
\newtheorem{corollary}[theorem]{Corollary}
\newtheorem{fact}[theorem]{Fact}
\theoremstyle{definition}
\newtheorem{definition}[theorem]{Definition}
\theoremstyle{remark}
\newtheorem{remark}[theorem]{Remark}

% Lean declaration names, directory and file names: typewriter text that may break after . / _ and, less readily,
% after : (never after the last character, so that a following punctuation mark stays on the line). The break points
% are explicit penalties, so they also work in statement heads, which amsthm builds in a box (the math-mode break
% points of \nolinkurl are lost there). \kern0pt after each character prevents ligatures and hyphenation. The
% arguments are plain names (letters, digits and . / _ : - ' only).
\makeatletter
\def\es@ttpen{0}
\def\es@ttcolpen{100}
\def\es@tt#1{{\upshape\ttfamily\let\es@ttpending\relax\expandafter\es@ttloop\detokenize{#1}\es@ttstop}}
\def\es@ttloop#1{\ifx\es@ttstop#1\else\es@ttpending\es@ttchar{#1}\expandafter\es@ttloop\fi}
\def\es@ttchar#1{#1\kern\z@
  \if.#1\def\es@ttpending{\penalty\es@ttpen\relax}\else
  \if/#1\def\es@ttpending{\penalty\es@ttpen\relax}\else
  \if_#1\def\es@ttpending{\penalty\es@ttpen\relax}\else
  \if:#1\def\es@ttpending{\penalty\es@ttcolpen\relax}\else
  \let\es@ttpending\relax\fi\fi\fi\fi}
\let\es@ttstop\relax
\DeclareRobustCommand{\lean}[1]{\ifmmode\text{\es@tt{#1}}\else\es@tt{#1}\fi}
\DeclareRobustCommand{\dirn}[1]{\ifmmode\text{\es@tt{#1}}\else\es@tt{#1}\fi}
% For the statement map (appendix): an underscore in roman text becomes a permitted line break, as in identifiers.
\newcommand{\esmapbreaks}{\protected\def\_{\ifmmode\nfss@text{\textunderscore}\else\textunderscore\allowbreak\fi}}
\makeatother
% The source of a statement: Lean name and file:line, printed after the statement head.
\newcommand{\leanref}[2]{Lean: \lean{#1}, \dirn{#2}}

% Geometry and combinatorics.
\newcommand{\orient}{\operatorname{orient}}
\newcommand{\conv}{\operatorname{conv}}
\newcommand{\hv}{\operatorname{hv}}
\newcommand{\peel}{\operatorname{peel}}
\newcommand{\sgn}{\operatorname{sgn}}
\newcommand{\R}{\mathbb R}
\newcommand{\Z}{\mathbb Z}
\newcommand{\N}{\mathbb N}
\newcommand{\C}{\mathbb C}
\newcommand{\abs}[1]{\lvert#1\rvert}
\newcommand{\card}[1]{\lvert#1\rvert}
